% Edition19: independently viewed French pp.60--63 and English pp.41--43.
% This is an audit fragment, not a replacement for the source's proof.
\egasection{§11：Grassmann 概形的法英回校}

\noindent\textbf{实际范围。}
本次逐段对读了法文手稿物理第60页“11. Grassmanniennes”起，经过第61、62页，至第63页末关于 $\operatorname{Grass}(E)$ 的手写补记；相应英文编辑稿为物理及印刷第41--43页。法文第59页和第64页仅用于核定前后边界，不计作本节回校正文。第63页中间虽已有一个被划去的旧§12标题，接在它下面的手写补记仍属§11；有效§12从法文第64页、英文第43页中段开始。现有§12及后文保持原位置，不能由“回校到§11”推断后文尚无译稿。

本次补回一处英译整段遗漏，区分多处英译错误与法文正确读法，并纠正一项旧译注的过度归责。实际查看的原页、哈希、段落映射及变更记录另存本轮工程记录。这里不把源页查看、编译或公式排版一致称为本节所有一般命题的独立证明。

\subsection*{一、恢复被英译漏掉的内在刻画条件}

法文第62页对相对维数 $m$ 的线性子簇作如下预告：它们是 $P$ 的闭子预概形，在 $S$ 上光滑，处处具有相对维数 $m$，并在每个 $s\in S$ 处具有“射影次数一”。英文第42页在这一句中漏掉“$P$ 的闭子预概形、在 $S$ 上光滑、处处相对维数 $m$”这些文字，只留下维数和次数的片段。旧中文因此采用过下述译注：
\begin{quote}
底本这句语法有缺口，只写出相对维数与逐纤维“射影次数一”，没有完整列明平坦性、有限表示等可能需要的条件。这里保留作者的预告，不将该残句单独表述为一个已经完备的充要判据。
\end{quote}
这项“缺口”判断只适用于所读英译，不能归给现在实际查看的法文。第十九版正文按法文补全以上三个条件，并以D19-11-09说明改动。作者随后明确说将在后章给出此刻画，且本段用不到它；这两句完整保留。本次没有补造原稿此处并不存在的证明，也没有声称已经审定任意拟凝聚 $E$ 情形下这项内在刻画的全部技术细节。

\Needspace{10\baselineskip}
\subsection*{二、原稿、英译与旧中文的错误归属}

\begin{longtable}{@{}p{.18\linewidth}p{.34\linewidth}p{.40\linewidth}@{}}
\toprule
对应新注 & 法文原页的实际读法 & 本次处理与旧注去向\\
\midrule
\endhead
D19-11-02 & 第60页已写 $\bigwedge^nG$ 及 $\bigwedge^nE_{S'}$。 & 英文第41页说明句漏外幂。沿用旧中文正确公式，明确D17-E04针对英译，不是法文。\\
D19-11-06 & 第61页正合列首项为 $E'$，下一行说明秩 $n$。 & 英文第41页误排首项 $E$。保留旧中文订正，明确其来源。\\
D19-11-07 & 第61页只出现一次“可由开浸入表示；另一方面”，正合列名为 $s$。 & 英文第42页重复短语并混写 $s,S$。去重内容不变；旧D17-E06关于选定分裂的条件继续保留。\\
D19-11-08 & 第62页手写插入“但不规定 $G$ 的秩”。 & 英译把它标为译者说明。正文恢复法稿已有的身份，不再称仅是英译者新增。\\
D19-11-10 & 第62页为 $(r-1)-1-m$，旁注相对余维 $m+1$。 & 英文多出的相邻加减号不归责于法稿；保留正确维数和空对偶边界。\\
D19-11-11 & 第63页所说的是 $L^0$ 的 $S'$-值点。 & 英文漏上标为 $L$；旧中文补上上标正确，现明确错误仅在英译。\\
\bottomrule
\end{longtable}

正文保留英文原脚注57--60：两处关于 \textit{sorites} 译法的询问、Cartan讨论班参考文献须更精确的询问，以及 paragraph 含义的询问。英文对 $G/Q$、末段手写字迹和向A.G.询问的编辑说明也按身份保留。没有据这些疑问虚构讨论班的年卷或原手稿的缺失句子。

\subsection*{三、有限表示下降：解释不等于订正原错}

法文第61页和英文第41页都说，有限表示的 $A$-模 $E$ 来自 $A$ 的某个有限型子环上的有限型模。旧中文曾把正文中的“有限型模”改成“有限表示模”，并写道：
\begin{quote}
底本这里只说“来自 $A$ 的某个有限型子环上的有限型模”。为使所述归约真正保持有限表示，取定义 $E$ 的有限表示矩阵的有限个系数所生成的子环，并在该子环上使用同一有限矩阵即可；正文将这项必要信息写明。这里不是断言任意有限型模都能如此下降其全部关系。
\end{quote}
有限矩阵这一解释本身正确，但不能据此把原稿措辞视为不足。若 $A_0$ 是由有限个系数生成的 $\mathbf Z$-子代数，则 $A_0$ 是诺特环，故其每个有限型模都是有限表示模。更直接地，用有限表示矩阵在 $A_0$ 上定义 $E_0$，便有 $E_0\otimes_{A_0}A\simeq E$。结合本节前面的有限型情形及与基变换相容性，正是原文所用的下降路线。

第十九版恢复法文的“有限型模”措辞，在D19-11-04解释以上理由；旧注的有限矩阵内容及“不能对任意有限型 $A$-模如此声称”的界限均保留在这里。更强的实际构造可以帮助解释原证明，但不应据此制造一项原文错误。

\subsection*{四、仍须保留的实际订正和有限范围核查}

\noindent\textbf{拟紧性。}
法文第61页与英文第41页均无条件声称 $\operatorname{Grass}_n(E)\to S$ 拟紧。这在一般拟凝聚 $E$ 下不成立。令 $S=\Spec k$、$E=\bigoplus_{i\geq0}ke_i$、$n=1$。对任意有限指标集 $F$，取 $j\notin F$；商映射 $e_j\mapsto1$、其余基向量映为0所给的点，不属于 $\bigcup_{i\in F}D_+(T_i)$。所以标准开覆盖不存在有限子覆盖。旧D17-E05的反例和正文有限型限制继续采用；本次确认这一错误也确实存在于法文，而非仅在英译。

\noindent\textbf{有限秩和空边界。}
法文在“局部自由”后划去“有限型”，但坐标开图的个数 $\binom rn$、维数 $n(r-n)$ 以及其后提升截面的论证，都在有限秩情形使用。正文保留这一明确范围；没有将无限秩仿射开图当作局部有限表示的光滑态射。$n>r$ 时没有秩 $n$ 商，故代表预概形为空；几何连通且非空的纤维断言只在 $n\leq r$ 部分使用。此处是说明适用范围，不声称已辨明作者划去该词的全部动机。

\noindent\textbf{坐标开图及向量丛方向。}
选定坐标分裂 $E=E'\oplus E''$ 且 $\operatorname{rank}E'=n$。在 $E'\to G$ 为同构的开图上，用此同构把商写成
\[
E'\oplus E''\longrightarrow E',\qquad (x,y)\longmapsto x+a(y),
\quad a:E''\longrightarrow E'.
\]
因 $E',E''$ 有限局部自由，$\mathscr Hom(E',E'')^\vee\simeq\mathscr Hom(E'',E')$；采用 $\mathbf V(M)=\Spec\operatorname{Sym}M$ 的约定，故参数预概形正是
\[
\mathbf V\bigl(\mathscr Hom(E',E'')\bigr).
\]
不能只看同态方向便把原式的两个模交换。若未选分裂，则所得到的是相应向量丛的挠子，不能无条件宣称有典范零截面。

同一有限坐标图说明 $k(s)$-值点可在 $s$ 的邻域上提升：有限个剩余域坐标分别在局部环中取代表，再缩小邻域。若 $k$ 为无限域，有限维仿射空间的每个非空开子集都有 $k$-有理点；与这些定义在 $k$ 上的坐标开图结合，得到原文的有理点结论。这两项核查只在上述有限秩情形进行。

\noindent\textbf{极性、包含和秩可变。}
若 $K=\ker(E\twoheadrightarrow G)$，则 $E^\vee/G^\vee\simeq K^\vee$，其秩为 $r-n$。从而对偶线性子簇的维数为 $r-n-1=r-m-2$，余维为 $m+1$。当 $m=r-1$ 时，对偶为空，按约定记维数 $-1$。

英文第43页对 $\mathbf P(G')\subset\mathbf P(G)$ 明确用充分条件“if”；法文第63页相应行末连词受扫描边缘影响。正文按可核的充分条件译出，不把旧中文的“等价于”冒作已辨出的法文原词。旧中文的逆向在这里仍可说明：设闭包含与到 $\mathbf P(E)$ 的映射相容，则把普遍商 $E\to\mathcal O(1)$ 限制到 $\mathbf P(G')$ 后，经由 $G$ 因子化；由伴随及 $p_*\mathcal O_{\mathbf P(G')}(1)\simeq G'$ 得到相容映射 $G\to G'$，它因 $E\to G'$ 满射而满射。秩零部分的结论直接成立。这个论证是译校说明，未代替法文连接词的实际辨读。

末段不指定秩的局部自由商，按局部有限秩理解；秩可在开闭部分变化。这样 $\coprod_{n\geq0}\operatorname{Grass}_n(E)$ 才代表所述函子。$n=0$ 对应唯一零商，$\operatorname{Grass}_0(E)=S$；不能把这个不交并解释成还包括无限秩商。

\subsection*{五、明确保留的限制}

法文第60页副题的划删字样、首段维数式中的字母，以及第63页被覆打的旧§12标题中间词仍未可靠辨定。英文相关形式与现有可辨内容已在译注中保留，没有用数学常识补成所谓作者终稿。第63页最后手写附记的左括号没有可见闭合符，中文只为排版闭合括号。

本节引用的Cartan讨论班证明未在此重新定位；一般拟凝聚模情形的Pl\"ucker闭浸入、以及作者预告的线性子簇内在刻画，均没有被本次有限范围核查升级为已经独立完整重证。已恢复全部可辨段落、公式及编辑脚注与“所有一般结论均审定”是两种不同结论。
